% 391_Puzzleteile_Kosmologie_De_ch.tex
% Johann Pascher, 9. Okt. 2026
% Lose Puzzleteile II: Kosmologie

\providecommand{\BM}{\textbf{[B]}}
\providecommand{\KM}{\textbf{[K]}}
\providecommand{\SM}{\textbf{[S]}}
\providecommand{\XM}{\textbf{[X]}}
\providecommand{\QM}{\textbf{[Q]}}
\providecommand{\EM}{\textbf{[E]}}
\providecommand{\SetzM}{\textbf{[SETZUNG]}}

\chapter*{Lose Puzzleteile II: Kosmologie}
\addcontentsline{toc}{chapter}{Lose Puzzleteile II: Kosmologie}

\begin{abstract}
Dok.~390 hat die Bausteine der Teilchenphysik eingeordnet, die sich nachträglich in der Literatur
fanden und ins Bild der FFGFT fügen. Dieses Dokument tut dasselbe für die Kosmologie. Es geht acht
Bausteine durch: Zwickys Ermüdung des Lichts, die Gravitationstheorie von Hoyle und Narlikar und
Wetterichs Universum ohne Expansion, Einsteins statisches Universum und die Steady-State-Theorie,
Diracs große Zahlen, die Hubble-Länge, Milgroms Beschleunigungsskala, das Problem der
Vakuumenergie und die kosmische Hintergrundstrahlung. Auch hier gilt die Reihenfolge: Die Bausteine
waren nicht Grundlage der FFGFT, sie fügen sich nachträglich ein. Anders als in der Teilchenphysik
liegt der kosmische Sektor aber außerhalb des Herleitungsanspruchs: $H_0$ ist an der Messung
kalibriert, der Exponent~10 ist eine Setzung, und das Standardmodell leitet diesen Sektor ebenso
wenig her. Die Puzzleteile zeigen deshalb keine Herleitungen, sondern Lesarten: Was die
Beobachtungen der Kosmologie tragen, lässt sich mit $\tilde T\cdot m=1$ ohne Expansion lesen, und
mehrere Stellen der Literatur haben das schon gezeigt. Zwei Bausteine, Vakuumenergie und
Hintergrundstrahlung, passen bisher nicht oder nur als Kandidat und werden als Grenze benannt.
\end{abstract}

% ============================================================
\section*{Statusmarkierungen}
% ============================================================
\begin{center}
\begin{tabular}{@{}lp{11cm}@{}}
\textbf{[SETZUNG]} & Motivierte Eingabe, keine Ableitung. \\[2pt]
\textbf{[B]} & Algebraisch bewiesen (Prüfskript, exakte Arithmetik). \\[2pt]
\textbf{[K]} & Numerisch bestätigt. \\[2pt]
\textbf{[S]} & Strukturell plausibel, Identifikation offen. \\[2pt]
\textbf{[E]} & Etablierte Mathematik oder Physik, übernommen. \\[2pt]
\textbf{[Q]} & Aus der Literatur übernommen, mit Quelle. \\[2pt]
\textbf{[X]} & Geprüft und negativ. \\
\end{tabular}
\end{center}
Die Marker geben den Status \emph{im Korpus} wieder; das Dokument leitet nichts neu her, sondern
ordnet ein. Prüfskript: \texttt{pruef\_391\_kosmologie.py}.

% ============================================================
\section{Reihenfolge und Grenze}
% ============================================================

\textbf{Zur Reihenfolge.} Wie in Dok.~390 war keiner der folgenden Bausteine Ausgangspunkt der FFGFT.
Die Theorie wurde aus $\tilde T\cdot m=1$ und der Geometrie des Trägers entwickelt; die
Literaturstellen sind erst nachträglich gefunden worden. Wo es heißt, ein Baustein füge sich an
einer Stelle des Korpus ein, ist das eine Übereinstimmung im Nachhinein, keine Herkunft.

\textbf{Zur Grenze.} Der kosmische Sektor wird im Korpus bewusst ausgeklammert. Das Standardmodell
leitet $H_0$ und $\Lambda$ ebenso wenig her, ein Vergleich wäre asymmetrisch (P39 in Dok.~190;
Dok.~365). Die Form der Beziehung zwischen Hubble-Länge und Planck-Länge liegt als
dimensionsloses $\xi$-Verhältnis fest (P20), der Exponent~10 in $H_0\propto\xi^{10}$ ist eine
Setzung \SetzM, und die Skala ist an der Messung von $H_0$ kalibriert (R137 in Dok.~190). Was
dieses Dokument einordnet, sind daher Lesarten der Beobachtungen, keine Vorhersagen ihrer Werte.

% ============================================================
\section{Rotverschiebung ohne Expansion}
% ============================================================

\subsection{Zwickys Ermüdung des Lichts (1929)}

\textbf{Was es enthielt.} Kurz nach Hubbles Entdeckung schlug Zwicky vor, die Rotverschiebung
ferner Galaxien entstehe nicht durch Expansion, sondern dadurch, dass Licht auf dem Weg Energie
verliert \QM~\cite{Zwicky1929}. Das war die erste statische Deutung der Hubble-Beziehung.

\textbf{Was fehlte.} Die Ermüdung des Lichts ändert nur das Photon. Ferne Uhren, etwa die
Lichtkurven von Supernovae, müssten dann im ursprünglichen Takt ablaufen. Beobachtet wird aber eine
Dehnung um $(1+z)$, und daran scheiterte die Idee.

\textbf{Wo es sich im Korpus einfügt.} Wegen $\tilde T\cdot m=1$ werden Wellenlänge und Uhrentakt
gemeinsam gebucht. Aus $1+z=e^{\xi_0 x}$ folgt die Zeitdehnung $(1+z)$ für alle $z$ \BM{} (R128 in
Dok.~190). Klassisches Tired Light scheitert, weil es nur die Photonseite ändert; die Massenlesart
kann diesen Fehler strukturell nicht machen (Dok.~312). Die Rotverschiebung ist achromatisch, die
Skala der Wegverlängerung ist über $H_0$ kalibriert (K6 in Dok.~190).

\subsection{Hoyle--Narlikar (1964) und Wetterichs Universum ohne Expansion (2013)}

\textbf{Was es enthielt.} Hoyle und Narlikar formulierten eine Gravitationstheorie, in der die
Teilchenmassen aus der Verteilung der Materie im Kosmos hervorgehen \QM~\cite{HoyleNarlikar1964}.
Wetterich zeigte 2013, dass ein expandierendes Universum mit konstanten Massen durch eine
Feldredefinition exakt in ein statisches Universum mit anwachsenden Teilchenmassen übergeht
\QM~\cite{Wetterich2013}. Beobachtbar sind nur dimensionslose Verhältnisse.

\textbf{Was fehlte.} Ein Grund, warum der Massenlauf die natürliche Lesart sein sollte und nicht
nur eine gleichwertige Umschreibung.

\textbf{Wo es sich im Korpus einfügt.} Dok.~312 führt die Äquivalenz von Expansion und Massenlauf
als \KM{} und verankert damit die Entartung aus Dok.~267: Rotverschiebung, Leuchtkraftentfernung,
baryonische akustische Oszillationen und Hintergrundtemperatur unterscheiden die beiden Lesarten
nicht. In der FFGFT ist der Massenlauf keine Zusatzannahme, weil Masse und Zeittakt über
$\tilde T\cdot m=1$ ohnehin dieselbe Größe sind. Dok.~312 hält die Grenze fest: Die Äquivalenz
zeigt, dass die statische Lesart möglich und beobachtungsgleich ist, nicht, dass sie richtig ist.

\subsection{Einsteins statisches Universum (1917) und die Steady-State-Theorie (1948)}

\textbf{Was es enthielt.} Einstein baute 1917 das erste relativistische Weltmodell als statischen,
räumlich geschlossenen Kosmos \QM~\cite{Einstein1917}. Bondi, Gold und Hoyle formulierten 1948 ein
Universum ohne Anfang, das im Großen immer gleich aussieht \QM~\cite{BondiGold1948,Hoyle1948}.

\textbf{Was fehlte.} Einsteins Modell brauchte eine eigens eingeführte kosmologische Konstante und
war instabil. Die Steady-State-Theorie brauchte eine fortlaufende Erzeugung von Materie, und die
Entdeckung der Hintergrundstrahlung 1965 galt als ihr Ende.

\textbf{Wo es sich im Korpus einfügt.} Die FFGFT lebt auf einem Träger ohne Rand,
$\partial T^4=\emptyset$ (Dok.~312). Dok.~313 rechnet die Frage nach dem Anfang auf dem kompakten
Zeitzyklus durch: Die beobachtbare Geschichte füllt einen Halbzyklus, und der
\enquote{Anfang} ist der Antipode des Zyklus, kein Ereignis \KM{} mit der Kalibrierung aus P20.
Weder eine kosmologische Konstante noch eine Materieerzeugung wird dafür gebraucht. Offen bleibt
nach Dok.~313, dass eine exakte Periodizität des Zyklus dynamisch nicht generisch ist und als
Randbedingung gesetzt werden muss.

% ============================================================
\section{Skalen}
% ============================================================

\subsection{Diracs große Zahlen (1937)}

\textbf{Was es enthielt.} Dirac bemerkte, dass das Verhältnis von elektrischer zu gravitativer
Kraft zwischen Elektron und Proton und das Weltalter in atomaren Zeiteinheiten beide etwa $10^{39}$
betragen, und vermutete dahinter einen Zusammenhang \QM~\cite{Dirac1937}.

\textbf{Was fehlte.} Ein Mechanismus. Diracs Folgerung, die Gravitationskonstante müsse mit der Zeit
abnehmen, ist durch Messungen ausgeschlossen.

\textbf{Wo es sich im Korpus einfügt.} Im Korpus ist die Beziehung zwischen Hubble-Länge und
Planck-Länge ein dimensionsloses $\xi$-Verhältnis (P20); der Exponent ist gesetzt, die Skala an
$H_0$ kalibriert. P39 hält fest, dass kein Rahmen diese Beziehung aus ersten Prinzipien herleitet.
Diracs Beobachtung erscheint damit als Verhältnis zweier Skalen derselben Geometrie, nicht als
zeitlicher Lauf einer Konstante. Diracs Name kommt im Korpus bisher nicht vor; die Einordnung ist
eine Lesart \SM.

\subsection{Die Hubble-Länge als Maximalskala}

\textbf{Was es enthielt.} Hubble fand 1929 die lineare Beziehung zwischen Entfernung und
Rotverschiebung \QM~\cite{Hubble1929}. Die Länge $c/H_0$ wird seither als Horizont eines
expandierenden Universums gelesen.

\textbf{Was fehlte.} Eine Lesart dieser Länge ohne Expansion.

\textbf{Wo es sich im Korpus einfügt.} Dok.~182 und Dok.~365 lesen $R_H=c/H_0$ als geometrische
Maximalskala des kompakten Trägers, nicht als Expansionshorizont. Ihr Wert ist an $H_0$ kalibriert.

\subsection{Milgroms Beschleunigungsskala (1983)}

\textbf{Was es enthielt.} Milgrom erklärte die flachen Rotationskurven der Galaxien ohne dunkle
Materie durch eine Beschleunigungsskala $a_0$, unterhalb derer das Newtonsche Gesetz abweicht
\QM~\cite{Milgrom1983}. Schon er bemerkte, dass $a_0$ in der Größenordnung von $cH_0$ liegt.

\textbf{Was fehlte.} Ein Grund für diese Nähe zur kosmischen Skala und eine Begründung der
Interpolationsformel.

\textbf{Wo es sich im Korpus einfügt.} Im Korpus ist $a_0=cH_0/(2\pi)$ \KM{} (Dok.~365); derselbe
Exponent~10, der $H_0$ trägt, reproduziert auch $a_0$, was bei einem reinen $H_0$-Fit nicht zu
erwarten wäre (Dok.~309, Querverankerung). Dok.~308 gewinnt die Form
$a=\sqrt{a_N^2+a_N a_0}$ mit $a_0=2a_{\mathrm{bg}}$ aus der Unruh-Kopplung \KM/\SM. Das Buch
\enquote{Dunkle Materie -- eine Illusion} führt die Rotationskurven aus. Der Exponent bleibt gesetzt.

% ============================================================
\section{Wo es noch nicht oder nur als Kandidat passt}
% ============================================================

\subsection{Vakuumenergie und kosmologische Konstante}

\textbf{Was es enthielt.} Die Nullpunktsenergie der Quantenfelder ist im Casimir-Effekt gemessen.
Als Quelle der kosmologischen Konstanten liegt ihre naive Abschätzung um viele Größenordnungen über
dem beobachteten Wert; Weinberg nannte das das Problem der kosmologischen Konstanten
\QM~\cite{Weinberg1989}.

\textbf{Stand im Korpus.} Die Verbindung zwischen Casimir- und Hintergrund-Energiedichte bei der
Länge $L_\xi$ ist eine Identität, kein Messbeleg \SM{} (R136 in Dok.~190). Die $\Lambda$-Formeln
älterer Dokumente treffen den beobachteten Wert nicht \XM{} (R137). In der statischen Lesart ist
$\Lambda$ ein Lesart-Artefakt ohne eigenen Referenten (Dok.~308) und wird nicht gesetzt
(Dok.~312); eine Herleitung des beobachteten Werts gibt es nicht.
Dieser Baustein fügt sich bisher nicht ein.

\subsection{Die kosmische Hintergrundstrahlung (1965)}

\textbf{Was es enthielt.} Penzias und Wilson fanden 1965 eine isotrope Strahlung von etwa
$3$~Kelvin \QM~\cite{PenziasWilson1965}. Ihr Spektrum und ihre Anisotropien gehören zu den
genauesten Daten der Kosmologie.

\textbf{Stand im Korpus.} Die Relationen für $T_{\mathrm{CMB}}$ aus $\xi$ ergeben eine reine Zahl,
die erst in der Einheit eV mit dem Messwert übereinstimmt; sie gelten als einheitenabhängige
Übereinstimmung, die Frage ist offen (R137, R70). Die Relation $T(z)$ in Dok.~039 ist durch
Messungen ausgeschlossen \XM. Der Faktor~3 vor der Peak-Folge ist eine Setzung, die
Peak-Selektion insgesamt offen (R139). In der Grundform stellt Dok.~388 den Kandidaten
$T_{\mathrm{CMB}}/m_e=(8\pi)^{1/4}\,\xi^{5/2}$ auf, der den FIRAS-Wert trifft \KM; der Vorfaktor
ist motiviert, aber nicht begründet \SM, und die Zufallsprüfung dort zeigt, dass ein Zufall nicht
ausgeschlossen ist. Gilt der Kandidat, verbindet er Temperatur und Hubble-Skala über
$T_{\mathrm{CMB}}^4=16\,H_0\,m_e^3$ \BM{} (Dok.~388), mit demselben gesetzten Exponenten~10.
Die Temperatur der Strahlung gehört zudem zu den Größen, die die Entartung aus Dok.~267 nicht
auflösen. Dieser Baustein fügt sich damit erst als Kandidat ein, nicht als gesicherte Stelle.

% ============================================================
\section{Was die Teile zeigen}
% ============================================================

In der Teilchenphysik (Dok.~390) fügten sich die Bausteine an Stellen ein, an denen der Korpus
etwas herleitet. In der Kosmologie fügen sie sich an Stellen ein, an denen der Korpus etwas
\emph{liest}. Zwicky, Hoyle und Narlikar, Wetterich, Einstein und die Steady-State-Theorie haben
jeweils ein Stück einer statischen Kosmologie geliefert; jedes Stück scheiterte für sich an einem
anderen Punkt, an der Zeitdehnung, an der fehlenden Begründung des Massenlaufs, an der Instabilität
oder an der Materieerzeugung. Mit $\tilde T\cdot m=1$ sind diese Punkte keine getrennten Probleme
mehr, weil Wellenlänge, Uhrentakt und Masse dieselbe Buchung sind.

Dirac und Milgrom haben je eine auffällige Nähe zwischen mikroskopischen und kosmischen Skalen
bemerkt. Im Korpus sind beide Nähen Verhältnisse derselben Geometrie, aber an einen gesetzten
Exponenten und eine kalibrierte Skala gebunden. Sie sind deshalb Querverankerungen, keine Belege.

Vakuumenergie und Hintergrundstrahlung passen bisher nicht oder nur als Kandidat. Das wird nicht
verdeckt: Die betreffenden Formeln sind im Register als \XM{} oder offen geführt, der
Temperatur-Kandidat aus Dok.~388 als \SM. Für den kosmischen Sektor gilt
damit dasselbe wie für schemaabhängige Größen in der Teilchenphysik: Ein Herleitungsanspruch wird
nicht erhoben, wo die Größen an der Messung oder an einer Konvention hängen.

% ============================================================
\section{Übersicht}
% ============================================================

\begin{center}
\small
\begin{tabular}{@{}>{\raggedright\arraybackslash}p{3.0cm}>{\raggedright\arraybackslash}p{1.7cm}>{\raggedright\arraybackslash}p{4.1cm}>{\raggedright\arraybackslash}p{3.8cm}@{}}
\toprule
Baustein & Jahr & Was fehlte & Im Korpus \\
\midrule
Zwicky, Ermüdung des Lichts & 1929 & keine Zeitdehnung & R128, K6; Dok.~312 \\
Hoyle--Narlikar, Wetterich & 1964, 2013 & Massenlauf nur gleichwertig, nicht begründet & Dok.~312, 267 \\
Einstein-Universum, Steady State & 1917, 1948 & Instabilität, Materieerzeugung & Dok.~313, 312 \\
Diracs große Zahlen & 1937 & kein Mechanismus & P20, P39 \SM \\
Hubble-Länge & 1929 & nur als Expansionshorizont gelesen & Dok.~182, 365 \\
Milgroms $a_0$ & 1983 & Nähe zu $cH_0$ unerklärt & Dok.~365, 309, 308 \\
Vakuumenergie, $\Lambda$ & 1989 & passt bisher nicht & R136, R137 \\
Hintergrundstrahlung & 1965 & nur Kandidat, Zufall nicht ausgeschlossen & Dok.~388; R137, R139, R70 \\
\bottomrule
\end{tabular}
\end{center}

% ============================================================
\section{Fazit}
% ============================================================

Auch in der Kosmologie lagen die Puzzleteile lange vor der FFGFT in der Literatur, und auch hier
waren sie nicht ihre Grundlage. Sechs der acht Bausteine fügen sich nachträglich ein, als Lesarten
der Beobachtungen ohne Expansion und als Querverankerungen zwischen mikroskopischen und kosmischen
Skalen. Zwei passen bisher nicht oder nur als Kandidat. Der Unterschied zur Teilchenphysik ist ehrlich zu benennen: Der
kosmische Sektor ist kalibriert und ausgeklammert, die Bausteine zeigen hier, dass die Theorie
die Kosmologie lesen kann, nicht, dass sie ihre Zahlen herleitet.

\vspace{1em}
\noindent\textit{Prüfskript:} \texttt{2/python/Dok391\_Skripte/pruef\_391\_kosmologie.py}
--- 51/51.

% ============================================================
\section*{Werkzeug und Methode}
\addcontentsline{toc}{section}{Werkzeug und Methode}
% ============================================================

Claude (Anthropic) wurde als Werkzeug eingesetzt für: Zusammenstellen der Fundstellen im Korpus,
Schreiben und Ausführen des Prüfskripts, Ausarbeiten der \LaTeX-Quellen nach den Vorgaben des
Autors. Das Prüfskript sichert jede Korpusangabe dieses Dokuments durch eine Textprüfung in der
genannten Kapiteldatei ab; schlägt eine Prüfung fehl, wird die Angabe nicht verwendet.

\begin{thebibliography}{99}

\bibitem{Zwicky1929}
F.~Zwicky, \textit{On the red shift of spectral lines through interstellar space},
Proc.\ Natl.\ Acad.\ Sci.\ USA \textbf{15} (1929) 773--779.

\bibitem{HoyleNarlikar1964}
F.~Hoyle, J.~V.~Narlikar, \textit{A new theory of gravitation},
Proc.\ R.\ Soc.\ Lond.\ A \textbf{282} (1964) 191--207.

\bibitem{Wetterich2013}
C.~Wetterich, \textit{A Universe without expansion},
Phys.\ Dark Univ.\ \textbf{2} (2013) 184--187; arXiv:1303.6878.

\bibitem{Einstein1917}
A.~Einstein, \textit{Kosmologische Betrachtungen zur allgemeinen Relativitätstheorie},
Sitzungsber.\ Preuß.\ Akad.\ Wiss.\ Berlin (1917) 142--152.

\bibitem{BondiGold1948}
H.~Bondi, T.~Gold, \textit{The steady-state theory of the expanding universe},
Mon.\ Not.\ R.\ Astron.\ Soc.\ \textbf{108} (1948) 252--270.

\bibitem{Hoyle1948}
F.~Hoyle, \textit{A new model for the expanding universe},
Mon.\ Not.\ R.\ Astron.\ Soc.\ \textbf{108} (1948) 372--382.

\bibitem{Dirac1937}
P.~A.~M.~Dirac, \textit{The cosmological constants}, Nature \textbf{139} (1937) 323.

\bibitem{Hubble1929}
E.~Hubble, \textit{A relation between distance and radial velocity among extra-galactic nebulae},
Proc.\ Natl.\ Acad.\ Sci.\ USA \textbf{15} (1929) 168--173.

\bibitem{Milgrom1983}
M.~Milgrom, \textit{A modification of the Newtonian dynamics as a possible alternative to the
hidden mass hypothesis}, Astrophys.\ J.\ \textbf{270} (1983) 365--370.

\bibitem{Weinberg1989}
S.~Weinberg, \textit{The cosmological constant problem}, Rev.\ Mod.\ Phys.\ \textbf{61} (1989) 1--23.

\bibitem{PenziasWilson1965}
A.~A.~Penzias, R.~W.~Wilson, \textit{A measurement of excess antenna temperature at 4080 Mc/s},
Astrophys.\ J.\ \textbf{142} (1965) 419--421.

\end{thebibliography}
